\section{楕円一ブロックの明示的可換復元}\label{er:sec:intro}

全体を通じて基礎体は$\CC$である。
三次中心正則元$g$を持つ次数付きAS-regular代数$S$から
$S[u_1,\ldots,u_m]/(g+P(u))$を作るとき、その残余圏を
可換三次超曲面から復元できるかを考える。
中心変数$u_i$は$S$の全ての元と可換とし、$P$は次数三、
\begin{equation}\label{er:eq:isolated}
 V(\partial_{u_1}P,\ldots,\partial_{u_m}P)=\{0\}\subset\mathbb A^m,
 \qquad m\ge2
\end{equation}
を仮定する。これは$P$のアフィン臨界点が原点のみであるという条件である。

使用する量子化は、$t^3\ne-1$に対する
\[
 f_t(x,y,z)=\frac{x^3+y^3+z^3}{3}+txyz
\]
のJacobian Poisson括弧の量子化である。
$h$を次数零の変形パラメータとし、$S_{h,t}$を
\begin{equation}\label{er:eq:relations}
\begin{aligned}
 [y,z]&=h\bigl(x^2+\tfrac t2(yz+zy)\bigr),\\
 [z,x]&=h\bigl(y^2+\tfrac t2(zx+xz)\bigr),\\
 [x,y]&=h\bigl(z^2+\tfrac t2(xy+yx)\bigr)
\end{aligned}
\end{equation}
で定める。$x,y,z,u_i$の内部次数は一である。
次の多項式と有理関数を用いる：
\begin{equation}\label{er:eq:coefficients}
\begin{aligned}
 b(h,t)&=1-\frac{h^2t^2}{4},&
 a(h,t)&=t+\frac{h^2}{3}\left(1+\frac{t^3}{4}\right),\\
 \vartheta(h,t)&=\frac{a(h,t)}{b(h,t)},&
 \Delta(h,t)&=a(h,t)^3+b(h,t)^3,\\
 U(h,t)&=1+\frac{h^2t^2}{2}+\frac{h^4t}{6}-\frac{h^4t^4}{48},\\
 W(h,t)&=\frac{(4-h^2t^2)\bigl(12t+h^2(t^3-8)\bigr)}{48}.
\end{aligned}
\end{equation}
$\mathscr U_t=\{h:U(h,t)b(h,t)\Delta(h,t)\ne0\}$は$0$を含む開集合である。
$\sym(xyz)$は六つの順序付き積の平均を表し、
\begin{equation}\label{er:eq:center}
 \Omega_{h,t}=\frac{U(h,t)}3(x^3+y^3+z^3)+W(h,t)\sym(xyz)
\end{equation}
と置く。中心性・正則性は第\ref{er:sec:algebra}節で確認する。

\begin{theorem}[一ブロックの明示的可換復元]\label{er:thm:main}
$t^3\ne-1$、$h\in\mathscr U_t$および\eqref{er:eq:isolated}を仮定する。
\[
 B_{h,t,P}=S_{h,t}[u_1,\ldots,u_m]/(\Omega_{h,t}+P(u)),
 \qquad
 Y_{h,t,P}=V\bigl(f_{\vartheta(h,t)}(x,y,z)+P(u)\bigr)
 \subset\PP^{m+2}
\]
とおく。$Y_{h,t,P}$は滑らかな三次超曲面であり、dg enhancementを保つ同値
\begin{equation}\label{er:eq:main-pair}
 (\mathcal K_{B_{h,t,P}},\rho_B)
 \simeq (\Ku(Y_{h,t,P}),\rho_Y)
\end{equation}
が存在する。ここで$\mathcal K_B$はOrlovの残余成分、
$\rho_B=\mathbf L_{\OO_B}\circ(1)$である。
従って$\mathcal K_B$はsmooth properである。
この可換模型は射影同型を除いて一意である。
楕円点曲線上の平行移動の有限位数性は仮定しない。
\end{theorem}

\begin{remark}[主張の範囲]\label{er:rem:scope}
定理は各複素パラメータにおける同値である。
全ての同値を底上の一つの代数的dg同値に選んだとは主張していない。
また$u_i$と$x,y,z$の本質的な混合項や、二つ以上の独立な無限位数楕円ブロックは対象外である。
smooth proper性はこの特定の残余成分について同値から得る結論であり、
任意のAS-regular中心三次商からの一般的帰結ではない。
\end{remark}

\subsection{先行研究と論理的な依存関係}
三次元Sklyanin代数の中心三次元と楕円曲線による商の表示は古典的である
\cite{SmithTate,Walton,DeLaetLeBruyn}。
Orlovの定理はGorenstein parameter零の特異圏を射影圏と同一視し、
正のparameterでは特異圏を残余成分として実現する\cite[Theorem 2.5]{Orlov}。
非可換行列因子分解との比較には\cite[Theorem 5.8]{CCKM}を用いる。
Etingof--Ginzburgの\cite[Theorem 3.6.2]{EG}にも楕円曲線とgraded factorizationの比較が現れる。
本稿では同結果のgenericという仮定を全パラメータへ読み替えず、
各$h\in\mathscr U_t$で中心商とOrlovの適用条件を確認する。

本稿の比較の要点は、楕円曲線の次数シフトに現れる平行移動を直線束のテンソルで共役し、
その三次周期性を保存したまま中心変数を追加することである。
有限群の双対同変化には\cite[Theorem 4.2, Section 8]{Elagin}を用いる。
可換模型の一意性は\cite[Theorem 1.7]{LZ}によるが、これは存在を与える定理ではない。
本稿の結果は既存の一楕円ブロック比較をHesse量子化へ適用し、
その変形方向と一定の可換模型を持つ代数的曲線を明示したものである。
これらの組合せについて先行研究上の初出を確定したとは主張しない。

\subsection{圏と次数の規約}
右加群を用いる。$\MF^\Gamma(A,w)$は有限射影$\Gamma$次数付き行列因子分解のdg圏の
pretriangulated hullを取り冪等完備化したものとする。
ここで行列因子分解自体は$D^2=w$を満たすcurved objectであり、通常の複体ではない。
一方、二対象間のHomには$d(\varphi)=D_Q\varphi-(-1)^{|\varphi|}\varphi D_P$が入り、
$w$の中心性から$d^2=0$となる。
$\Gamma=\ZZ$、$\deg w=3$なら、内部シフトは自然な周期性
\begin{equation}\label{er:eq:period}
 (1)^3=(3)\simeq[2]
\end{equation}
を備える。以下の同値では、この周期性も記録する。

\section{中心三次元と楕円商}\label{er:sec:algebra}
\begin{proposition}\label{er:prop:center}
$h\in\mathscr U_t$では$S_{h,t}$はnoetherianな三次元Koszul AS-regular domainで、
Hilbert級数は$(1-z)^{-3}$である。
\eqref{er:eq:center}は中心正則元であり、$\Omega_{0,t}=f_t$。
さらに滑らかな三次曲線
\[
 E_{h,t}=V\bigl(x^3+y^3+z^3+3\vartheta(h,t)xyz\bigr)
\]
と$L=\OO_{E_{h,t}}(1)$について
\begin{equation}\label{er:eq:thcr}
 S_{h,t}/(\Omega_{h,t})\simeq B(E_{h,t},L,t_{p_{h,t}}),
 \qquad
 p_{h,t}=[1-ht/2:-1-ht/2:-h]
\end{equation}
が成り立つ。原点は$o=[1:-1:0]$である。$h=0$では$t_{p_{0,t}}=\Id$と解釈する。
\end{proposition}
\begin{proof}
$q=ht/2$、$\alpha=1-q$、$\beta=-1-q$、$\gamma=-h$とすると、関係式は
$\alpha yz+\beta zy+\gamma x^2=0$とその巡回置換である。
$h\ne0$では$\alpha\beta\gamma=h b(h,t)\ne0$かつ
\[
 (\alpha^3+\beta^3+\gamma^3)^3-(3\alpha\beta\gamma)^3
 =-27h^3\Delta(h,t)\ne0.
\]
従って楕円型Sklyanin代数の既知の定理を適用できる\cite{SmithTate,Walton}。
その点曲線は
$\alpha\beta\gamma(x^3+y^3+z^3)-
(\alpha^3+\beta^3+\gamma^3)xyz=0$である。
$\alpha^3+\beta^3+\gamma^3=-3h a(h,t)$を代入すると$E_{h,t}$を得る。
$1+\vartheta^3=\Delta/b^3\ne0$なので滑らかである。

中心元の規格化を確認するため、$S_3=x^3+y^3+z^3$、
$C=xyz+yzx+zxy$、$D=yxz+xzy+zyx$と書く。
\cite[Section 2, p.~3]{DeLaetLeBruyn}の中心三次元の表示は
\[
 G=\alpha(\beta^3-\gamma^3)C+
   \beta(\gamma^3-\alpha^3)D+
   \gamma(\alpha^3-\beta^3)S_3.
\]
関係式を$x,y,z$倍して足すと$C-D=hS_3+q(C+D)$である。
この式を用いて直接整理すると
\begin{equation}\label{er:eq:central-identity}
 G=-9h\left(\frac U3 S_3+\frac W6(C+D)\right)=-9h\Omega_{h,t}.
\end{equation}
従って$h\ne0$では$\Omega_{h,t}$は標準中心元の非零スカラー倍であり、
中心正則で、その商が\eqref{er:eq:thcr}となる。
$h=0$では代数は多項式環、$U(0,t)=1$、$W(0,t)=t$であるから
$\Omega_{0,t}=f_t$である。この場合は滑らかな平面三次曲線の通常の斉次座標環であり、
同じ商の記述が成立する。
\end{proof}

\begin{remark}
$h=0$における$S_{0,t}$自体のpoint schemeは$\PP^2$であって、$E_{0,t}$ではない。
$E_{0,t}$は中心三次元で割った商の射影曲線である。
この二つを混同すると零パラメータを含む族の記述を誤る。
また中心元の公式の証明では\cite{DeLaetLeBruyn}のSection 2の表示を使う。
そのarXiv v1の冒頭のTheorem 1に印刷された表示とSection 2の表示には係数の差があるため、
\eqref{er:eq:central-identity}を別途多項式として検算する。
\end{remark}

\section{楕円曲線上のrotationの共役}\label{er:sec:rotation}
この節では$E$を滑らかな平面三次曲線、$L=\OO_E(1)$、$\sigma=t_p$とし、
$S$をnoetherianな三次元二次AS-regular代数、$g\in Z(S)_3$を正則元で
$S/(g)=B(E,L,\sigma)$となるものとする。

\begin{lemma}\label{er:lem:block}
$\MF^\ZZ(S,g)$は$\Perf(E)$とdg Morita同値である。
この同値の下で内部シフトは
\[
 \rho_p=T_{\OO_E}\circ\sigma_*\circ(-\otimes L)
       =t_{-p}^*\rho_0,
 \qquad \rho_0=T_{\OO_E}\circ(-\otimes L)
\]
に対応する。$T$はspherical twistである。
\end{lemma}
\begin{proof}
三次元AS-regular代数のparameterは三であり、三次中心正則元で割ると
Gorenstein parameterは零になる。
\cite[Theorem 5.8]{CCKM}のcokernel関手および
\cite[Theorem 2.5(iii)]{Orlov}のwindowを合成すると
\[
 \MF^\ZZ(S,g)\simeq \Dsg^{\mathrm{gr}}(S/(g))
 \simeq D^b(\qgr(S/(g)))\simeq\Perf(E)
\]
を得る。ここではcokernelとwindowのdg模型を使う。
同じ有限射影分解・dg quotientによってHom複体の比較が得られるので、
三角圏の抽象的同値だけを用いているわけではない。
THCRの右加群の規約では、tailsの次数シフトは
$\sigma_*(-\otimes L)$である。
windowを一つ移動させる際の差は$\OO_E$の評価射のconeであり、
従って残余側の次数シフトは$T_{\OO_E}\sigma_*(-\otimes L)$になる。
最後に、平行移動は$\OO_E$を保つのでそのpullbackは$T_{\OO_E}$と可換である。
\end{proof}

\begin{lemma}[次数零直線束による共役]\label{er:lem:conjugacy}
$P_q=\OO_E(q-o)$、$A_q=(-\otimes P_q)$と置くと
\begin{equation}\label{er:eq:conjugacy}
 A_q\rho_0A_q^{-1}\simeq t_q^*\rho_0.
\end{equation}
特に$\rho_p=A_{-p}\rho_0A_{-p}^{-1}$であり、$A_p$は$(\Perf(E),\rho_p)$から
$(\Perf(E),\rho_0)$への同値を与える。$p$の位数に条件はない。
\end{lemma}
\begin{proof}
テンソル関手との共役により$A_qT_{\OO_E}A_q^{-1}=T_{P_q}$である。
sourceを第一因子、targetを第二因子とするFourier--Mukai規約で、
$T_{P_q}$のkernelは
\[
 (P_q^{-1}\boxtimes P_q)\otimes I_\Delta[1]
\]
である。$E\times E$の二つの方向のfiberへの制限を計算すると、see-sawの原理により
\[
 (P_q^{-1}\boxtimes P_q)\otimes\OO(-\Delta)
 \simeq\OO(-\Gamma_q),\qquad
 \Gamma_q=\{(x,y):y=x-q\}
\]
となる。実際、第一座標$x$を固定すると両辺の次数$-1$の直線束の類は
$-x+q$であり、第二座標$y$を固定すると類は$-y-q$である。
よって$T_{P_q}=t_q^*T_{\OO_E}$。
$A_q$は$(-\otimes L)$と可換なので\eqref{er:eq:conjugacy}を得る。
\end{proof}

\begin{lemma}[周期性を含むブロック比較]\label{er:lem:periodic-block}
$E=V(f_E)$ならば、内部シフトと\eqref{er:eq:period}を保つdg Morita同値
\[
 \MF^\ZZ(S,g)\simeq\MF^\ZZ(\CC[x,y,z],f_E)
\]
が存在する。
\end{lemma}
\begin{proof}
両側を補題\ref{er:lem:block}により$\Perf(E)$と同一視し、
補題\ref{er:lem:conjugacy}の$A_p$を用いる。
これにより内部シフトをintertwineする自然同型を得る。
二つの三次周期性をこの同型に沿って比較した差は、同値関手の可逆な自然自己変換である。
$\HH^0(E)=H^0(E,\OO_E)=\CC$だから、この差は非零スカラーである。
その三乗根を用いてintertwinerを再規格化すれば、三次周期性も一致する。
負次数のHochschild cohomologyは零であり、既にfactorization側で与えられた
内部シフトのcoherenceをこの同値で移す。
従って単に自己同値の同型類を合わせるだけでなく、次節の有限群作用に必要な
周期性のデータも保存できる。
\end{proof}

\section{中心変数の追加と有限次数差の降下}\label{er:sec:suspension}
この節は、ブロックの比較を全非可換cubicへ移すための生成性を示す。
$T=\CC[u_1,\ldots,u_m]$、$A=S\otimes T$、$w=g+P$とする。
次数群を
\[
 \Gamma=(\ZZ e_0\oplus\ZZ e_1)/(3e_0-3e_1)
\]
とし、$S_d$の次数を$de_0$、$T_d$の次数を$de_1$とする。

\begin{lemma}[中心変数に対する外部積生成]\label{er:lem:TS}
$S$が非負次数付きnoetherian AS-regular代数、$g\in Z(S)_3$が正則で、
$P$が\eqref{er:eq:isolated}を満たせば、外部積はdg Morita同値
\begin{equation}\label{er:eq:TS}
 \Perf\bigl(\MF^\ZZ(S,g)\otimes_\CC\MF^\ZZ(T,P)\bigr)
 \simeq\MF^\Gamma(A,w)
\end{equation}
を与える。第一・第二因子の内部シフトはそれぞれ$(e_0)$、$(e_1)$に対応する。
\end{lemma}
\begin{proof}
外部積の微分は$D_Q\otimes1+(-1)^{\deg}1\otimes D_K$であり、その二乗は$g+P$である。
有限自由factorization間のHomを行列成分ごとに計算すると、右辺のHom複体は左辺の
二つのHom複体のtensor totalizationに一致する。
非負の内部次数と固定された有限個の生成元のshiftにより、各Hom複体は下に有界であり、
各total次数では和は有限である。これが完全忠実性を与える。

本質的全射性には$u$方向の微分を使う。
任意のfactorization$F$の微分を$D$とすると、$S$と$u_i$が独立なので
$\partial_{u_i}$は$S[u]$の導分であり、
\[
 D(\partial_{u_i}D)+(\partial_{u_i}D)D
   =(\partial_{u_i}P)\Id_F.
\]
従ってJacobianイデアル$J(P)$は$F$上で零ホモトピックに作用する。
\eqref{er:eq:isolated}より、十分大きい$d$に対し全ての$u_i^d$は$J(P)$に属する。
$u_i^d$による乗法のconeを順に取ったKoszul対象は$F$のshiftを直和因子に持つ。
特異圏でその対象を有界複体として表すと、そのcohomologyは
$u_1^d,\ldots,u_m^d$で殺される。
$(u_1,\ldots,u_m)$の冪による有限filtrationの各商は
$C=S/(g)=A/(w,u)$上の有限生成次数付き加群である。
従って右辺は、$C$-加群のrestrictionをstabilizeした対象で生成される。

ここでrestrictionが特異圏に降りることを確認する。
$u_1,\ldots,u_m$は$A/(w)$上の中心正則列である。
実際、途中で一部の$u$を零にしても$g+$残りの三次式は
$S$次数によるfiltrationで正則な最高記号$g$を持つ。
一つの$u_i$についても、先に$u_i=0$とした商でこの元が正則であることから
$u_i$が非零因子となる。
従って$C$は$A/(w)$上有限射影次元を持ち、restrictionはperfectをperfectへ送る。

$C$-加群$M$の十分高いsyzygyを取れば、その$S$上の射影次元は高々一となる。
実際、自由$C$-加群の$S$上の射影次元は一なので、syzygyを一つ取るごとに
次数二以上の$\Ext_S$は一つdimension shiftし、$S$の有限大域次元により最後には消える。
その二項$S$-射影分解と$g$の零ホモトピーは
$Q\in\MF^\ZZ(S,g)$を与える。
$M$を$A$上で分解するには、この二項分解と$u_1,\ldots,u_m$のKoszul分解をtensorすればよい。
$w$の零ホモトピーには、前者の$g$の零ホモトピーと
\[
 P=\sum_i u_i\,\frac{\partial_{u_i}P}{3}
\]
から得るKoszulホモトピーを足す。
この有界分解を周期化すると、restrictionされた$M$のstabilizationは
\[
 Q\boxtimes K_P,\qquad
 K_P=\left(\bigwedge^\bullet\CC^m\otimes T,
       \sum_i u_i\iota_i+\frac{\partial_{u_i}P}{3}\,\varepsilon_i\right)
\]
に一致する。$\varepsilon_i$は外積、$\iota_i$は縮約で、$K_P$の微分の二乗は$P$である。
二項分解の完全性、Koszul分解の完全性、およびこの式は、周期化のcohomologyを
直接同定し、余分な有限射影尾部は特異圏で零となる。

最後に、$C$は$\ZZ e_0$に次数を持つ一方、
$\Gamma/\ZZ e_0\simeq\ZZ/3$である。
任意の$\Gamma$次数付き$C$-加群をこの三剰余類に分ければ、
上の$\ZZ$次数付き議論と有限個のshiftですべての対象が得られる。
従って粗化次数でのみ現れる加群を落とさず、外部積が全生成する。
\end{proof}

\begin{lemma}[通常次数への降下]\label{er:lem:descent}
補題\ref{er:lem:periodic-block}のように第一因子の内部シフトと三次周期性を保つdg同値は、
\eqref{er:eq:TS}を通じ、通常の整数次数について
\[
 \MF^\ZZ(S[u],g+P)\simeq
 \MF^\ZZ(\CC[x,y,z,u],f_E+P)
\]
を誘導する。この同値は通常次数の$(1)$を保つ。
\end{lemma}
\begin{proof}
$\Gamma\to\ZZ$を$e_0,e_1\mapsto1$で定めると、核は$\ZZ/3$である。
その双対群$H\simeq\mu_3$は、整数次数の代数において$u$だけを三乗根倍する。
$P$の次数が三なのでcurvatureは不変である。
$\Gamma$次数付きfactorizationは、整数次数付きfactorizationの$H$-同変対象と同一視される。

その同変圏には双対群$H^\vee\simeq\ZZ/3$が文字をtensorすることで作用する。
外部積の表示では生成元は$\rho_T\rho_S^{-1}$であり、その三乗の自明化は
二つの$(3)\simeq[2]$の比較である。
第一因子の同値がこのデータを保つので、$H^\vee$-同変化にも同値を誘導する。
有限群の位数が可逆であることと冪等完備性を用い、
双対同変化で元の圏を回復する定理\cite[Theorem 4.2]{Elagin}を適用する。
同論文Section 8のdg模型、あるいは有限群のaveraging idempotentで、この比較はdg同値として行える。
$e_0$のシフトは全過程で保存されるため、通常次数の$(1)$も保存される。
\end{proof}

\subsection{主定理の証明}
\begin{proof}[定理\ref{er:thm:main}の証明]
命題\ref{er:prop:center}により$(S,g)=(S_{h,t},\Omega_{h,t})$は
補題\ref{er:lem:block}の全ての仮定を満たす。
$E=V(f_{\vartheta})$なので、周期性を保つブロック比較と中心変数の追加・降下から
\[
 \MF^\ZZ(S_{h,t}[u],\Omega_{h,t}+P)
 \simeq\MF^\ZZ(\CC[x,y,z,u],f_{\vartheta}+P)
\]
を得る。$S_{h,t}[u]$は次元$m+3$のnoetherian AS-regular代数であり、
$\Omega_{h,t}+P$は中心正則である。
その商のGorenstein parameterは$m$であり、
Orlov windowはこのMF圏を
$\langle\OO_B,\ldots,\OO_B(m-1)\rangle^\perp$へ実現する。
degree shiftはその残余成分の$\mathbf L_{\OO_B}(1)$に移る。
可換側も同じwindowを用いるので\eqref{er:eq:main-pair}を得る。

可換方程式の全偏微分が零なら、まず\eqref{er:eq:isolated}により$u=0$である。
残る$f_{\vartheta}$は$1+\vartheta^3\ne0$だからアフィン臨界点は原点のみである。
従って$Y$は射影的に滑らかである。
$\Ku(Y)$は滑らかな射影多様体のadmissible成分なのでsmooth properである。
最後に、別の滑らかな可換cubicのpairと同値ならば、二つの可換pairのenhanced同値を得る。
これはadmissible成分のFourier--Mukai同値として表され、
\cite[Theorem 1.7]{LZ}から二つのcubicは同型になる。
次元$m+1\ge3$ではその同型は超平面束を保ち、射影同型である。
\end{proof}

\begin{corollary}[任意の滑らかな三変数三次式]\label{er:cor:ternary}
滑らかな三変数三次式$g$を
$g(x)=\lambda f_t(Mx)$、$\lambda\ne0$、$M\in\mathrm{GL}_3(\CC)$とHesse形にする。
$g$のJacobian括弧を量子化し、$h'=\lambda\det(M)h$、
$\Omega_h^g(x)=\lambda\Omega_{h',t}(Mx)$で中心元を正規化する。
このとき$h'\in\mathscr U_t$では、
$S_h^g[u]/(\Omega_h^g+P)$の可換pair模型は
\[
 V\bigl(\lambda f_{\vartheta(h',t)}(Mx)+P(u)\bigr)
\]
である。さらに$Mx$を中心変数の一次形式だけ平行移動した場合も、同じ座標置換で結論を移せる。
\end{corollary}
\begin{proof}
滑らかな平面三次曲線のHesse標準形は\cite[Lemma 1]{ArtebaniDolgachev}による。
$y=Mx$の下で定数三重ベクトルは$\det(M)$倍となり、対称化は線形変換と可換なので
$S_h^g\simeq S_{h',t}$である。
商の方程式を$\lambda$で割り、定理\ref{er:thm:main}を$P/\lambda$に適用する。
最後に方程式を$\lambda$倍して元の座標に戻す。
中心線形シフトは次数付き代数の同型であり、内部シフトも保つ。
\end{proof}

\section{復元されたcubicの二次から始まる変形}\label{er:sec:second-order}
次の主張は、定理\ref{er:thm:main}で明示した可換cubicの族自体に関するものである。
非可換族の全体との相対同値を仮定しない。

\begin{theorem}[markedな二次変形の非自明性]\label{er:thm:second-order}
$t$と$P$を固定する。$Y_{h,t,P}$の方程式は
\begin{equation}\label{er:eq:second}
 f_t+P+\frac{h^2(1+t^3)}{3(1-h^2t^2/4)}xyz
\end{equation}
である。従って$h$に関して一次のmarked変形は自明だが、
$\CC[h]/(h^3)$上のmarked変形は自明でない。
$q=h^2$をパラメータとする族のKodaira--Spencer類は
\[
 \frac{1+t^3}{3}[xyz]\in
 J(f_t+P)_3,
 \qquad J(F)=\CC[x,y,z,u]/(\partial F),
\]
であり、これは非零である。
\end{theorem}
\begin{proof}
\eqref{er:eq:coefficients}から直接
\[
 \vartheta(h,t)-t=
 \frac{h^2(1+t^3)}{3(1-h^2t^2/4)}
\]
を得る。滑らかな射影超曲面のembeddedな一次変形を射影座標変換と方程式のスカラー倍で割った空間は
Jacobian環の三次部分である。Euler等式により$F$自身もJacobianイデアルに含まれる。
次元三以上では、この記述は$H^1(T_Y)$による抽象変形の接空間とも一致する。

$f_t$の三つの偏微分は$x^2+tyz,y^2+tzx,z^2+txy$であり、正則列をなす。
従ってそのJacobian環はHilbert級数$(1+z)^3$を持ち、三次部分の次元は一である。
同環で
\[
 x^2y=-t y^2z=t^2xz^2=-t^3x^2y
\]
なので、$1+t^3\ne0$から$x^2y=0$となり、同様に全ての指数型$(2,1,0)$の単項式が零となる。
また$x^3=y^3=z^3=-txyz$である。
従って三次部分は$xyz$で生成され、次元一であるから$[xyz]\ne0$。
$t=0$でも同じ議論が成立し、$xyz$は三次の唯一の平方自由単項式である。
変数が分離しているので$J(f_t+P)=J(f_t)\otimes J(P)$であり、
$[xyz]\otimes1$は非零のままである。

一次の変形項が零なので、$h^2$での座標変換が消せるのはJacobianイデアルの三次部分だけである。
$h$で始まる座標変換が一次を保つ場合、その一次項は$F$の無限小射影自己同型である。
滑らかな三次超曲面では、偏微分の一次係数syzygyが存在しないこととEuler等式により
それはスカラーEuler方向に限られ、射影的には零である。
よって非零の上記Jacobian類は$\CC[h]/(h^3)$上で消せない。
\end{proof}

\begin{remark}
この結論はmarked変形、または局所moduli stackについてのものとして述べている。
特別なcubicの有限自己同型による商を取ったcoarse moduliの局所座標で、
同じ二次の次数がそのまま最初の非零次数になるとは主張しない。
また、ファイバーごとの\eqref{er:eq:main-pair}のみから、非可換族の相対Kodaira--Spencer類を
上の式と同一視してはいない。
\end{remark}

\section{同じ可換模型を持つ非可換な代数的曲線}\label{er:sec:constant-model}
固定した$t$では復元されたcubicが動く一方、$t$も動かすと正反対の現象が現れる。

\begin{theorem}[ファイバーごとに同じpairを持つ量子化族]\label{er:thm:isocategorical}
$\tau^3\ne-1$と\eqref{er:eq:isolated}を満たす$P$を固定し、
$Y_0=V(f_\tau+P)$とする。
$(h,t)=(0,\tau)$を通る平面曲線
\begin{equation}\label{er:eq:curve}
 \mathscr C_\tau:
 h^2t^3+3\tau h^2t^2+12t+4h^2-12\tau=0
\end{equation}
のその点を含む適切な滑らかな開部分上に、次数別平坦な代数的族
$B_{h,t,P}$がある。次を満たす。
\begin{enumerate}
 \item $h$への射は$(0,\tau)$でétaleであり、全ての複素ファイバーで
 $(\mathcal K_{B_{h,t,P}},\rho_B)\simeq(\Ku(Y_0),\rho_{Y_0})$。
 \item $B_{h,t,P}$の$(0,\tau)$における次数付き結合代数の一次変形は非自明である。
 特にこの族は可換代数の定数族ではない。
 \item この曲線の関数体上で、Sklyaninの平行移動点$p_{h,t}$は無限位数である。
\end{enumerate}
ここで主張しているのはファイバーごとの同値であり、
一つの相対同値による自明化ではない。
\end{theorem}
\begin{proof}
$\vartheta(h,t)=\tau$の分母を払うと\eqref{er:eq:curve}になる。
その$t$による偏微分は$(0,\tau)$で12である。
従ってこの点を含む唯一の既約成分を取り、適切な開近傍に制限すると、
曲線は滑らかで$h$への射はétaleとなる。
さらに$t^3+1$、$U$、$b$、$\Delta$が非零となるよう縮める。
いずれも$(0,\tau)$で非零である。
ここでは$\vartheta=\tau$なので、定理\ref{er:thm:main}の可換模型は常に$Y_0$である。

代数の表示はこの曲線上の正則関数を係数とする有限表示である。
各幾何ファイバーの$S$のHilbert級数は$(1-z)^{-3}$で一定である。
滑らかな曲線上の有限生成加群のtorsionはその台でfiber次元を増加させるので、
各次数成分は局所自由になる。
中心元による各次数の乗法は全fiberで単射であり、その余核も局所自由である。
従って商$B$は次数別平坦である。

局所環の完備化では$t=\tau+O(h^2)$である。
従って$\CC[h]/(h^2)$上で
\[
 [x,y]=h(z^2+\tau xy).
\]
右辺の係数は$B_0=\CC[x,y,z,u]/(f_\tau+P)$で零でない。
実際、唯一の定義方程式は三次であり、非零の二次式を消さない。
可換な定数変形へ同型な代数は可換なので、この一次変形は自明ではない。

最後に、点曲線は一定の$E_\tau$であり、$p_{h,t}$はその上の点である。
原点$o$の局所パラメータ$z/x$で
\[
 (z/x)(p_{h,t})=-\frac h{1-ht/2}=-h+O(h^2)
\]
となる。標数零の楕円曲線の形式群では、任意の正整数$r$について
$[r](v)=rv+O(v^2)$である。
従って$[r]p$の局所パラメータは$-rh+O(h^2)$であり、恒等的に零ではない。
よって$p$は関数体上で無限位数である。
\end{proof}

\begin{corollary}[固定したHesse係数でのgenericな位数]\label{er:cor:nontorsion}
固定した$t^3\ne-1$に対しても、$p_{h,t}$は$\CC(h)$上で無限位数である。
各固定整数$r\ge1$に対し、$[r]p_{h,t}=o$となる$h\in\mathscr U_t$は有限個である。
従って非常に一般的な$h$で$p_{h,t}$は無限位数である。
\end{corollary}
\begin{proof}
滑らかな楕円曲線族の零切断に沿う形式群にも$[r](v)=rv+O(v^2)$が成り立つので、
前の証明と同じ一次係数が非torsion性を与える。
二つの切断$[r]p$と$o$の一致集合は底の閉集合であり、恒等ではない。
一次元の有限型既約底の真の閉集合は有限である。
全$r$の例外集合の和は可算なので、最後の主張が従う。
\end{proof}

\section{相対中心商と、相対復元に残る一段階}\label{er:sec:relative}
この節では$t$を固定し、
$R=\CC[h,(U(h,t)b(h,t)\Delta(h,t))^{-1}]$とする。
$\mathscr E\to\Spec R$は$f_{\vartheta}=0$で定まる滑らかな楕円曲線族、
$\mathcal L=\OO_{\mathscr E}(1)$、$p=p_{h,t}$、$\sigma=t_p$とする。
$S_R$と$\Omega_R$は式\eqref{er:eq:relations}、\eqref{er:eq:center}で定義する。

\begin{proposition}[中心商の相対THCR表示]\label{er:prop:relative-THCR}
$\mathcal L_d=\mathcal L\otimes\sigma^*\mathcal L\otimes\cdots\otimes(\sigma^{d-1})^*\mathcal L$
と置くと、自然な次数付き同型
\begin{equation}\label{er:eq:relative-THCR}
 S_R/(\Omega_R)\simeq
 R\oplus\bigoplus_{d\ge1}H^0(\mathscr E,\mathcal L_d)
\end{equation}
が存在する。右辺の積は$\sigma$で捩った切断の積である。
この同型と両辺の各次数成分の形成は、任意の可換$R$-代数への底変更と両立する。
\end{proposition}
\begin{proof}
各$S_{R,d}$は有限生成$R$-加群であり、全閉fiberにおける次元は一定である。
$R$は$\CC[h]$の局所化であるからPIDであり、torsion成分があればその台でfiber次元が増える。
従って各次数は有限自由である。
同様に$\Omega_R$の乗法の余核も全fiberでHilbert級数
$(1-z^3)/(1-z)^3$を持ち、次数別有限自由である。

$d\ge1$では各fiberの$\mathcal L_d$の次数は$3d>0$である。
楕円曲線の$H^1$の消滅とcohomology and base changeにより、
$H^0(\mathscr E,\mathcal L_d)$は階数$3d$の有限局所自由$R$-加群で、その形成は任意の底変更と可換である。
$d=0$についても$\pi_*\OO_{\mathscr E}=\OO_{\Spec R}$で同じ主張が成立する。

三つの座標切断から$S_R$の右辺への写像を作る。
二次関係式が満たされることは各複素fiberで点曲線のgraphによる表示から従う。
右辺の各次数が$R$-torsion-freeなので、この恒等式は$R$上でも成立する。
同じ理由で$\Omega_R$の像は零である。
得られた商からの写像は$h\ne0$では古典的THCR表示、$h=0$では平面三次曲線の
projective normalityによる通常のsection ring表示であり、全fiberで同型である。
各次数の有限自由加群間の写像にNakayamaの補題を適用すると\eqref{er:eq:relative-THCR}を得る。
各次数が自由であり、切断の形成も底変更と可換なので、最後の主張も従う。
\end{proof}

\begin{remark}[証明できた相対命題と残る問題]\label{er:rem:relative-gap}
命題\ref{er:prop:relative-THCR}は、非可換三次元ブロックの中心商を実在する代数的底上で記述している。
しかし、それだけで
$\MF_R(S_R,\Omega_R)$と$\Perf(\mathscr E/R)$を
任意の底変更と両立して同一視したわけではない。
さらに$P(u)$を加えた残余圏について、対角のperfect性まで得たわけでもない。

正則な底についてのOrlov比較の候補はBurke--Stevenson
\cite[Definition 5.1, Theorem 6.4]{BurkeStevenson}である。
同定理には、全体の代数の左右noetherian性、有限injective dimension、
$\RHom_A(A_0,A)$のGorenstein型、$\chi$条件が必要である。
ファイバーごとのAS-regular性だけを、これらの相対的な条件の証明として扱ってはならない。
本稿はこの経路で相対同値とその底変更compatibilityを全て確認したとは主張しない。

明示的な滑らかな族$Y\to\Spec R$の相対Kuznetsov成分はsmooth properな可換模型として存在する。
これと実際の非可換量子化の相対残余圏を同定することが次の課題である。
一般の中心混合項および複数の無限位数ブロックの全生成は、さらに別の問題として残る。
\end{remark}